%\documentclass{...}

\input{./def/yf-formatting}

\usepackage{amsfonts, amsmath, amssymb, amsthm}
\usepackage{comment}
\usepackage{graphicx}
\usepackage{ifthen}
\usepackage{latexsym}
%\usepackage{times}
\usepackage[normalem]{ulem}

\input{./def/yf-def}

\def\extraspacing{\vspace{5mm} \noindent}

\def\best{\mathit{best}}
\def\size{\mathit{size}}

%\newboolean{solver}\setboolean{solver}{true}
\newboolean{solver}\setboolean{solver}{false}
\ifthenelse{\boolean{solver}}{\includecomment{sol}}{\excludecomment{sol}}

\pagenumbering{gobble}

\begin{document}

.

\vspace{50mm}

\begin{center}
    \uline{SC3020-CE4301-CZ4031: Quiz 3 (Make-Up)}
\end{center}

\vspace{20mm}
\begin{center}
    Name:
    \uline{\hspace{30mm}}
    \hspace{10mm}
    Student ID:
    \uline{\hspace{30mm}}
\end{center}

\pagebreak

\extraspacing {\bf Problem 1 (40 marks).} Consider three transactions:
\myitems{
    \item $T_1:$ R(A); R(B); W(C); \vspace{-2mm}
    \item $T_2$: R(B); W(A); W(D); \vspace{-2mm}
    \item $T_3$: R(C); R(A); W(B); R(A);
}
where R($X$) means reading $X$ from the disk, and W($X$) means writing $X$ to the disk. Suppose that all the transactions are executed with the same isolation level. For each schedule below, indicate the strongest isolation level (among \ttt{READ UNCOMMITTED}, \ttt{READ COMMITTED}, \ttt{REPEATABLE READS}, and \ttt{SERIALIZABLE}) that can produce the schedule. You need to justify your answer.

\vgap

(a) R1(A); R1(B); R2(B); W1(C); Commit1; R3(C); R3(A); W2(A); W2(D); Commit2; W3(B); R3(A); Commit3;

Remark: the numbers in a schedule denote transaction IDs; e.g., R1($A$) represents an R($A$) operation from $T_1$.

\vgap

(b) R1(A); R1(B); R2(B); W2(A); W1(C); R3(C); W2(D); Commit2; R3(A); W3(B); Commit1; R3(A); Commit3;

\begin{sol}
    \extraspacing {\bf Solution.} (a)
    \myitems{
        \item Unrepeatable read: Yes ($T_3$ reads A before and after W2(A)).
        \item Dirty read: No
    }
    We can now conclude that the strongest isolation level is \ttt{READ COMMITTED}.

    %\extraspacing {\bf Solution.}
    %(b) By swapping non-conflicting operations, it is possible to convert the schedule to the serial schedule $T_1, T_2, T_3$. Therefore, the strongest isolation level is \ttt{SERIALIZABLE}.

    (b) Due to W1(C); R3(C), it is \ttt{READ UNCOMMITTED}.
\end{sol}

\pagebreak


%\vspace{100mm}

\extraspacing {\bf Problem 2 (40 marks).} Let $t$ be a tuple from a single-attribute relation. The current history of $t$ is $(10, 100, \text{NULL})$, and the transaction with ID 100 has committed. Give a schedule satisfying all the requirements below:

\myenums{
    \item It involves two transactions, whose IDs are 101 and 102, respectively.

    \item It can be produced by the MVCC protocol.

    \item Executing the schedule will create the following history of $t$:
    \begin{center}
        (10, 100, 101), (20, 101, 102), (30, 102, NULL).
    \end{center}

    \item It contains non-repeatable reads.
}

\noindent Your schedule must contain these instructions: READ$(t)$, LOCK$(t)$, WRITE$(t, \text{value})$, and COMMIT.

\noindent Remark: You get 20 marks if your schedule satisfies all requirements except requirement 4.


\begin{sol}
\extraspacing {\bf Solution.}

\begin{center}
    \begin{tabular}{c|l|l}
        & \(T_{1}\) (ID 101) & \(T_{2}\) (ID 102) \\
        \hline
        1 & & \(\text{READ}(t)\) \\
        2 & \(\text{LOCK}(t)\) & \\
        3 & \(\text{WRITE}(t, 20)\) & \\
        4 & COMMIT & \\
        5 & & \(\text{READ}(t)\) \\
        6 & & \(\text{LOCK}(t)\) \\
        7 & & \(\text{WRITE}(t, 30)\) \\
        8 & & \(\text{COMMIT}\)
    \end{tabular}
\end{center}

\end{sol}
\pagebreak

\extraspacing {\bf Problem 3 (20 marks).} Suppose that the database system sees the following log after a reboot.

\begin{center}
\begin{tabular}{l}
$\langle \text{START } T_{1} \rangle$ \\
$\langle \text{START } T_{2} \rangle$ \\
$\langle T_{1}, A, 100, 200 \rangle$ \\
$\langle T_{2}, B, 50, 150 \rangle$ \\
$\langle \text{COMMIT } T_{1} \rangle$ \\
$\langle \text{START } T_{3} \rangle$ \\
$\langle T_{3}, C, 300, 250 \rangle$ \\
$\langle T_{2}, D, 80, 120 \rangle$ \\
$\langle \text{COMMIT } T_{2} \rangle$ \\
$\langle T_{3}, A, 200, 350 \rangle$ \\
$\langle \text{START } T_{4} \rangle$ \\
$\langle T_{4}, B, 150, 400 \rangle$
\end{tabular}
\end{center}

\noindent Describe the steps taken by the recovery procedure.

\begin{sol}
\extraspacing {\bf Solution.} $T_1$ and $T_2$ had committed before the crash while $T_3$ and $T_4$ had not. The redo phase processes the log forward and performs three disk modifications (in this order):
\myitems{
    \item set $A$ to 200;
    \item set $B$ to 150;
    \item set $D$ to 120.
}
The undo phase processes log backward and performs two disk modifications (in this order):
\myitems{
    \item set $B$ to 150;
    \item set $A$ to 200;
    \item set $C$ to 300.
}
\end{sol}

% \pagebreak
% .
% \vspace{100mm}


\end{document}
